\section{A variational lower bound on the ground energy}
\label{sec:hastings_variational_lower_bound}

% Hastings, arXiv:cond-mat/0508554v3, proj3new.tex, lines 856--877 (eq. varnum)
% and 1005--1009.

For a translation-invariant Hamiltonian $H=\sum_iH_i$ on $N$ sites, whose terms
$H_i$ are translates of one another, \cite{Hastings2006Solving} observes that the
ground energy per site is at least the lowest eigenvalue of $H_i$, and that the
same holds after adding to $H_i$ any combination of its iterated commutators
with $H$, because these commutators add up to zero over the translates. Write
$\operatorname{ad}_H(X)=HX-XH$, and $\lambda_{\min}(K)$ for the lowest eigenvalue
of a Hermitian $K$, the largest $\lambda$ with $K-\lambda\ge0$.

\begin{lemma}[Iterated commutators of Hermitian operators]
    \label{thm:hastings_commutator_iter_hermitian}
    \lean{Matrix.isHermitian_smul_commutatorIter,
        Matrix.conjTranspose_commutatorIter, Matrix.commutatorIter}
    \leanok
    For Hermitian $H$ and $A$, a real number $a$ and $k\ge0$, the operator
    $a\,i^k\operatorname{ad}_H^k(A)$ is Hermitian.
\end{lemma}
\begin{proof}\leanok
    By induction, $(\operatorname{ad}_H^kA)^\dagger=(-1)^k\operatorname{ad}_H^k(A)$,
    and $\overline{i^k}(-1)^k=i^k$.
\end{proof}

\begin{lemma}[The variational lower bound for unitary conjugates]
    \label{thm:hastings_variational_lower_bound}
    \lean{Matrix.card_mul_le_of_mulVec_eq_smul_of_conj,
        Matrix.card_mul_sInf_spectrum_le_of_conj,
        Matrix.posSemidef_sub_card_smul_of_conj, Matrix.commutatorIter_conj,
        Matrix.commutatorIter_sum, Matrix.commutatorIter_self_succ}
    \leanok
    \uses{thm:hastings_commutator_iter_hermitian}
    Let $H_0$ be an operator, let $U_1,\dots,U_N$ be unitaries, and let
    $H=\sum_iU_iH_0U_i^\dagger$ with every $U_i$ commuting with $H$. For real
    constants $a_1,\dots,a_m$ let
    \begin{align}
        K=H_0+\sum_{k=1}^ma_k\,i^k\operatorname{ad}_H^k(H_0).
    \end{align}
    If $K-\lambda\ge0$, then $H-N\lambda\ge0$; in particular every eigenvalue
    of $H$, and so its ground energy $E_0$, is at least $N\lambda$. This
    generalizes equation (varnum) of \cite{Hastings2006Solving}, whose terms are
    the translates of one term; the source's form is
    Lemma~\ref{thm:hastings_variational_lower_bound_sites}.
\end{lemma}
\begin{proof}\leanok
    Conjugation by a unitary $U$ that commutes with $H$ commutes with
    $\operatorname{ad}_H$, so
    $\sum_iU_i\operatorname{ad}_H^k(H_0)U_i^\dagger=\operatorname{ad}_H^k(H)=0$
    for $k\ge1$, and $\sum_iU_iKU_i^\dagger=H$. Since $U_i(K-\lambda)U_i^\dagger\ge0$
    and $U_iU_i^\dagger=1$, summing gives $H-N\lambda\ge0$. For an eigenvector
    $\Omega$ of $H$ with eigenvalue $E$,
    $0\le\langle\Omega,(H-N\lambda)\Omega\rangle=(E-N\lambda)\|\Omega\|^2$.
\end{proof}

\begin{lemma}[The variational lower bound]
    \label{thm:hastings_variational_lower_bound_sites}
    \lean{Matrix.card_mul_le_of_mulVec_eq_smul_of_sitePerm,
        Matrix.card_mul_sInf_spectrum_le_of_sitePerm}
    \leanok
    \uses{thm:hastings_variational_lower_bound,thm:hastings_site_invariant_sum}
    Let a finite group $G$ act on a finite set of sites, such as the
    translations of a periodic lattice of any dimension or a group of
    automorphisms of a finite graph, and let $P_g$ be the unitary that moves the
    configurations of the sites along $g$. For the sum
    $H=\sum_{g\in G}P_gH_0P_g^\dagger$ of the translates of one term $H_0$, and
    $K$ as in Lemma~\ref{thm:hastings_variational_lower_bound}, every
    eigenvalue of $H$ is at least $|G|\lambda$ whenever $K-\lambda\ge0$. For
    Hermitian $H_0$, with $\lambda=\lambda_{\min}(K)$ and $N=|G|$,
    \begin{align}
        E_0/N\ge\lambda_{\min}\Bigl(H_0+\sum_{k=1}^ma_k\,i^k
        \operatorname{ad}_H^k(H_0)\Bigr).
    \end{align}
    Here $N=|G|$ counts the summands. When $G$ acts simply transitively on the
    sites, as the translations of a periodic lattice do, $N$ is the number of
    sites. This is equation (varnum) of \cite{Hastings2006Solving}, lines
    856--877, for translation-invariant $H=\sum_iH_i$; there the terms are
    written with $[H_i,H]=-\operatorname{ad}_H(H_i)$, which only changes the
    signs of the arbitrary constants.
\end{lemma}
\begin{proof}\leanok
    \uses{def:hastings_site_permutation}
    The map $g\mapsto P_g$ is multiplicative, so conjugation by $P_h$ permutes
    the translates and fixes $H$, and each $P_h$ is a unitary commuting with $H$.
    Lemma~\ref{thm:hastings_variational_lower_bound} applies. For Hermitian
    $K$, its real spectrum is finite, and every spectral value is at least
    its infimum. The spectral theorem therefore gives
    $K-\lambda_{\min}(K)\ge0$, yielding the stated specialization.
\end{proof}

The Gaussian expansion discussed at lines 1005--1009 of
\cite{Hastings2006Solving} corresponds, in the convention used here, to
$a_k=0$ for odd $k$ and
$a_k=[(k-1)!!/k!](q/\Delta E^2)^{k/2}$ for even $k$. The printed variance
power requires correction, and $i^k$ already supplies the alternating sign
\cite{gap:hastings05_gaussian_series_coefficients}. The norm-convergent Gaussian
identity is proved in QICLean; the finite-polynomial theorem above allows
arbitrary real coefficients.

\begin{lemma}[The Gaussian commutator expansion]
    \label{lem:hastings_gaussian_commutator_series}
    \lean{TNLean.Lattice.gaussianFilteredTerm_eq_exp_commutator,
        TNLean.Lattice.hasSum_gaussianFilteredTerm_commutator}
    \notready
    \uses{def:hastings_gaussian_filtered_term}
    Let $H$ be Hermitian, let $W$ be any operator, and let $q\ge0$ and
    $\Delta E>0$. The actual Gaussian-filtered term of variance
    $h=q/\Delta E^2$ satisfies
    \begin{align}
        \widetilde W
         & =\exp\bigl(-\tfrac h2\operatorname{ad}_H^2\bigr)(W) \\
         & =\sum_{r=0}^{\infty}\frac{(-h/2)^r}{r!}
        \operatorname{ad}_H^{2r}(W),
    \end{align}
    where $\operatorname{ad}_H(W)=HW-WH$ and the series converges in
    operator norm. In particular, the variance has power $r$ at order
    $2r$, as required by the correction of the printed coefficients at
    lines 1005--1009 of \cite{Hastings2006Solving}, documented in
    \cite{gap:hastings05_gaussian_series_coefficients}.
\end{lemma}
\begin{proof}
    \notready
    Apply the Gaussian exponential and norm-convergent series identities
    at variance $q/\Delta E^2$. Nonnegativity of $q$ identifies that
    variance with the nonnegative variance of the actual filter.
\end{proof}

\begin{lemma}[The Gaussian-filtered variational lower bound]
    \label{thm:hastings_gaussian_variational_lower_bound}
    \lean{TNLean.Lattice.sum_conj_gaussianFilteredTerm,
        TNLean.Lattice.card_mul_sInf_spectrum_gaussianFilteredTerm_le_of_conj,
        TNLean.Lattice.card_mul_sInf_spectrum_gaussianFilteredTerm_le_of_sitePerm}
    \leanok
    \uses{thm:hastings_variational_lower_bound, def:hastings_gaussian_filtered_term}
    Let $H_0$ be Hermitian, let $U_1,\ldots,U_N$ be unitaries, and suppose
    each $U_i$ commutes with $H=\sum_i U_iH_0U_i^\dagger$.
    Let $K$ be the Gaussian average of $e^{itH}H_0e^{-itH}$, with variance
    $q/\Delta E^2$, where $q\ge0$ and $\Delta E>0$. Then
    $H=\sum_iU_iKU_i^\dagger$, and every eigenvalue $E$ of $H$ satisfies
    $E\ge N\lambda_{\min}(K)$. For a simply transitive translation action,
    $N$ is the number of sites, so $E_0/N\ge\lambda_{\min}(K)$.
    This is the filtered choice discussed in \cite{Hastings2006Solving},
    source lines 1005--1009, in the variational bound (varnum).
\end{lemma}
\begin{proof}\leanok
    \uses{thm:hastings_variational_lower_bound, thm:hastings_gaussian_filtered_spectral,
        thm:hastings_filtered_term_translate}
    Covariance of the Gaussian average gives
    $U_iKU_i^\dagger=\widetilde{U_iH_0U_i^\dagger}$, because $U_i$ fixes $H$.
    The filtered terms sum to $H$. Apply
    Lemma~\ref{thm:hastings_variational_lower_bound} with local term $K$
    and all commutator coefficients zero.
\end{proof}
